\begin{table}[H]
\raggedright
\caption{ENMO distribution and stationarity checks.}
\label{tab:sup-stationarity}
\footnotesize
\begingroup
\setlength{\tabcolsep}{4pt}
\renewcommand{\arraystretch}{1.08}
\noindent\textit{Panel A.} Skewness and kurtosis of candidate activity transformations.\par\vspace{2pt}
\begin{tabular*}{\textwidth}{@{\extracolsep{\fill}}lcc@{}}
\toprule
Measure & Skewness & Kurtosis \\
\midrule
ENMO (raw) & 1.321 & 5.403 \\
ENMO (sqrt) & 0.444 & 2.974 \\
ENMO (log) & -0.504 & 3.432 \\
\bottomrule
\end{tabular*}
\par\vspace{6pt}\noindent\textit{Panel B.} Share of individual series flagged by linear trend and KPSS diagnostics.\par\vspace{2pt}
\begin{tabular*}{\textwidth}{@{\extracolsep{\fill}}lcc@{}}
\toprule
Variable & Prop. linear trend & Prop. KPSS non-stationary \\
\midrule
Activity & 0.088 & 0.118 \\
Fatigue & 0.235 & 0.206 \\
Pain & 0.353 & 0.176 \\
Stress & 0.265 & 0.176 \\
\bottomrule
\end{tabular*}
\par\vspace{3pt}\noindent\parbox{\textwidth}{\footnotesize\textit{Note.} ENMO was analysed on the log scale because the transformed distribution was closest to symmetry. KPSS = Kwiatkowski-Phillips-Schmidt-Shin test.}
\endgroup
\end{table}
